\chapter{Yogic Practices for Health: An Introduction}
\label{ch:practices-intro}

\emph{What are the families of yoga practice, what does each one do, and why are they done in a particular order?}
Before the practical chapters teach each technique step by step, the syllabus asks the candidate to know the practices in theory: their meaning, their purpose and their traditionally credited benefits. This chapter is that map. It covers yoga's view of health, the cleansing actions, the loosening practices and surya namaskara, asana, and pranayama and dhyana. The methods are in the practical chapters, referred to at each point.

\begin{unverified}
This chapter is written from notes blocks whose fact check is still pending. Most points follow the official guidebook (its second, fifth, seventh, eighth and ninth chapters) and the Common Yoga Protocol. A reviewer must check it before print.
\end{unverified}

\section{Health and the practices}

\subsection*{What health means}
The guidebook quotes the World Health Organization as describing health as complete physical, mental, social and spiritual well-being, and not only the absence of disease. The WHO Constitution of 1948 itself names physical, mental and social well-being; ``spiritual'' is the guidebook's addition, so state it as the guidebook's wording. The Sanskrit word for health, \emph{svāsthya}, comes from \emph{sva}, self, and \emph{stha}, abiding: health is abiding in one's own self.

In yoga's picture of a healthy body, the muscles are relaxed, the joints are loose and save energy, energy is used efficiently, and the organ systems stay coordinated even under stress. The guidebook notes that postures such as naukasana, utkatasana and bhujangasana build strength in the way push-ups, lunges and squats do, that vrikshasana trains balance on one leg, and that yoga moves the joints in ways that walking and running do not, which widens their range.

The texts have their own account. The Shvetashvatara Upanishad names lightness of body, health, freedom from craving, a clear complexion, a pleasant voice, an agreeable odour and scanty excretions as the first signs of progress in yoga, and the Hatha Pradipika echoes these as marks of success. The Yoga Sutra speaks of beauty, grace, strength and firmness as the perfection of the body.

\subsection*{The families of practice}
The yoga practices used for health and well-being are the cleansing actions (kriya), the loosening and flowing movements, asana, pranayama, mudra and bandha, meditation, relaxation and prayer. As traditionally taught, kriya cleanses; asana steadies and strengthens; pranayama calms and lightens; meditation steadies the mind; relaxation relieves tension. The Common Yoga Protocol of the Ministry of AYUSH is a standard short set of these practices for everyone.

The practices can also be grouped by the level they work on. For the body: asana, pranayama, shavasana for rest, the cleansing actions and the locks. The protocol says that the cleansing actions remove accumulated waste and that bandhas and mudras, linked with pranayama, help control the mind. For the mind: meditation, prayer and mantra. For the emotions: devotion, bhakti yoga. Overall, a teacher chooses a balanced routine to suit the person.

\subsection*{Principles of practice}
Practise with awareness, within your capacity, regularly and without strain; begin with a prayer and end with peace; and let the teacher guide the choice of practices. The five-point guide to a yogic life gives the same picture in one line: proper exercise, proper breathing, proper relaxation, proper diet, and positive thinking with meditation (see \cref{ch:guiding-principles}).

These practices are for prevention and well-being. They do not replace medical treatment.

\section{Shatkarma: the six cleansing actions}
\label{ch:shatkarma-intro}

\emph{Shat} means six and \emph{karma} (or \emph{kriyā}) means action: the shatkarmas are the six cleansing practices of the body. The guidebook gives their purpose: to clear the internal organs of waste that the body does not clear well by itself, to keep the three humours (vata, pitta, kapha) in balance, and to make the body ready for pranayama and the higher practices. Describe them as traditionally used for cleansing, never as a cure.

The six, and what each cleans:
\begin{itemize}
\item \textbf{Dhauti}: the digestive tract up to the stomach (vaman with water, danda with a tube, vastra with a cloth).
\item \textbf{Basti}: the colon (jala basti with water, sthala basti with air).
\item \textbf{Neti}: the nasal passage (jala neti with salt water, sutra neti with a thread or tube).
\item \textbf{Trataka}: the eyes, by steady gazing at one point without blinking.
\item \textbf{Nauli}: the abdomen, by churning the abdominal muscles.
\item \textbf{Kapalabhati}: the front of the head and the lungs, by quick, forceful out-breaths.
\end{itemize}
This is the order of the text; the order differs between texts and schools.

The Hatha Pradipika says the cleansing actions are for those with excess fat or phlegm, and that pranayama is enough for others. After jala neti, kapalabhati is done to clear any water left in the nose. Dhauti, basti and nauli are never taught in a beginners' class; they are learned from a qualified teacher. For the exam, know dhauti and neti and practise kapalabhati (see \cref{ch:cleansing}).

\paragraph{Facts often asked.} The Gheranda Samhita calls kapalabhati \emph{bhālabhāti} and gives it three forms, vatakrama, vyutkrama and sheetkrama. \emph{Achamana} is the ritual sipping of water; \emph{abhyanga} is oil massage of the body. \emph{Piṇḍa} and \emph{brahmāṇḍa} are the individual and the universe, the microcosm and the macrocosm.

\section{The loosening practices and surya namaskara}
\label{ch:vyayama-intro}

Movement keeps the joints supple and their fluid moving. Daily life uses many joints very little; yoga gives them gentle, full movement, and balances strength with stretch. Three groups of practice do this.

\begin{itemize}
\item \textbf{Sukshma vyayama} (micro-circulation practice): small, gentle movements joined with breath and attention, worked from the top of the body down: the neck, then the shoulders, elbows, wrists and fingers, then the trunk, the hips, and the knees, ankles and feet. Each exercise has three features: the breath, the point of attention, and the movement. It loosens, warms and relaxes, and needs no special preparation and little time.
\item \textbf{Sthula vyayama} (macro-circulation practice): larger, stronger movements for the whole body. The syllabus names two, sarvanga pushti and hrid gati (engine daud).
\item \textbf{Surya namaskara}: a flowing sequence of twelve positions joined with the breath, which bends the spine forward and back and warms the body for asana. It joins movement, breath and intention.
\end{itemize}
Keep the three side by side: sukshma is small and joint by joint; sthula is large and whole-body; surya namaskara is a flowing sequence. The subtle and the gross work complement each other, and together they give balanced fitness. They are traditionally credited with flexibility, better circulation, deeper breathing, healthy body weight, and a body and mind ready for asana, pranayama and meditation.

The guidebook gives only a short account of sukshma vyayama and credits its modern form to one teaching line; other schools describe it differently. The methods are in the practical chapters (see \cref{ch:vyayama} and \cref{ch:surya-namaskara}).

\section{Yogasana: meaning, principles and benefits}
\label{ch:asana-intro}

\subsection*{Meaning}
An asana is a way of sitting or holding the body. The Yoga Sutra says that a posture should be steady and comfortable: \emph{sthira-sukham āsanam}. The next sutra adds the relaxation of effort, so the three qualities of a good posture are stability, comfort and effortlessness. The guidebook says that asanas bring poise, balance and strength and calm the mind as a preparation for meditation, and that the physical gains are incidental to that calm.

\subsection*{Principles of practice}
Move slowly with the breath; hold the final position steadily and with awareness; do not strain or compare yourself with others; keep within your own capacity; finish by relaxing. Practise with the bowels and bladder empty and the stomach light, in a ventilated place on a level floor with a mat, and set unwanted thoughts aside. Do not practise during an illness, or after surgery, a fracture or a sprain, until a physician allows it. The workshop also teaches alignment, breath awareness, balance, stability and mindfulness as guiding principles; that is a teaching list, not a Board list.

\subsection*{The five movements of the spine}
\begin{itemize}
\item \textbf{Flexion}: bending forward, as in padahastasana.
\item \textbf{Extension}: bending back, as in bhujangasana and ushtrasana.
\item \textbf{Lateral flexion}: bending to the side, as in trikonasana.
\item \textbf{Rotation}: twisting, as in ardha matsyendrasana.
\item \textbf{Axial extension}: lengthening upward, as in tadasana.
\end{itemize}

\subsection*{Benefits}
The Hatha Pradipika names three benefits of asana: health (\emph{ārogya}), steadiness (\emph{sthairya}) and lightness of the body (\emph{aṅgalāghava}). The guidebook adds calm, flexibility, range of movement, strength, balance, and a higher level of fitness with less risk of injury. It also makes therapeutic claims, which this book does not repeat. Simhasana, the lion posture, is traditionally taught as helpful for the throat, ears and nose. All of these are for prevention and well-being, not cure. The postures of the syllabus are in \cref{ch:asanas}.

\section{Pranayama and dhyana: meaning and benefits}
\label{ch:pranayama-intro}

\subsection*{Pranayama}
\emph{Prāṇa} is the life force and \emph{āyāma} is extension or control, so pranayama is the control of prana, worked through the breath. The Yoga Sutra stresses the breath as a whole (inhalation, exhalation, retention); hatha yoga stresses retention (kumbhaka). A breath has three parts: \emph{pūraka}, the in-breath; \emph{kumbhaka}, retention, inside or outside; and \emph{recaka}, the out-breath.

The seats named for pranayama and meditation are siddhasana, padmasana and ardha padmasana, svastikasana, sukhasana and vajrasana. Before starting, check: eyes closed; trunk straight, with the head, neck and back in line; shoulders and belly relaxed; hands on the knees in a mudra; body still; mood calm.

Pranayama can be joined with postures and with meditation. It is traditionally credited with a calmer mind, a steadier breath, better lung capacity and endurance, less anxiety and stress, and a mind prepared for meditation. Breath holding and forceful breathing are not for beginners; teach the practices that need no holding first (see \cref{ch:pranayama}).

\subsection*{Dhyana}
Dhyana is meditation: the unbroken flow of attention towards one object, the seventh of Patanjali's limbs. The three inner limbs follow each other. Dharana is the effort to hold the mind on one object and to bring it back again and again; when that hold becomes unbroken, it is dhyana; when even the sense of the meditator fades and only the object shines, it is samadhi.

There are two broad ways to practise: concentration, fixing the mind on a breath, a mantra or a flame; and mindfulness, watching thoughts without judging them. The two can be combined. A beginner keeps a fixed time and a quiet place, sits a few minutes at first and then longer, on a comfortable seat with the back erect, and watches the breath; if sleepy, takes a deep breath and breathes out firmly.

The guidebook names calm and self-realisation as the aims of meditation. Meditation is traditionally credited with calm, clarity, better concentration and less reaction to stress. Call it a supportive practice for well-being, not a treatment. The method is in \cref{ch:meditation}.

\begin{recall}
\begin{itemize}
\item Svasthya: abiding in one's own self.
\item Shatkarma: dhauti, basti, neti, trataka, nauli, kapalabhati; for cleansing, not cure.
\item Sukshma (small, joint by joint), sthula (large, whole body), surya namaskara (flowing sequence).
\item Asana: \emph{sthira-sukham}; stable, comfortable, effortless. Hatha Pradipika: arogya, sthairya, angalaghava.
\item Pranayama: puraka, kumbhaka, rechaka. Dharana, then dhyana, then samadhi.
\end{itemize}
\end{recall}

\subsection*{Worked examples}
\begin{itemize}
\item \textbf{Model answer: what are the shatkarmas and why are they done?} Six cleansing actions: dhauti, basti, neti, trataka, nauli and kapalabhati. They clear waste the body does not clear well by itself, keep the three humours in balance, and prepare the body for pranayama. They are learned from a qualified teacher.
\item \textbf{Model answer: distinguish dharana, dhyana and samadhi.} Dharana is holding the mind on one object with effort; dhyana is when that attention flows without a break; samadhi is when only the object shines and the sense of a separate meditator fades.
\item \textbf{A student asks whether yoga will cure their back pain.} Explain that yoga supports well-being and may help them feel better, but does not replace medical care; ask them to see their doctor and to tell the teacher what the doctor advises.
\end{itemize}

\subsection*{Check yourself}
\begin{enumerate}
\item What does svasthya mean?
\item Do yoga practices replace medical treatment?
\item What does shatkarma mean, and which one cleans the eyes?
\item What should follow jala neti?
\item What does the Gheranda Samhita call kapalabhati?
\item In what order is sukshma vyayama done?
\item Which is the subtle exercise and which the gross?
\item How does the Yoga Sutra describe asana?
\item Name the five movements of the spine.
\item Which three benefits of asana does the Hatha Pradipika name?
\item Name the three parts of pranayama.
\item What is the difference between dharana and dhyana?
\end{enumerate}

\subsection*{Answers}
\begin{enumerate}
\item Abiding in one's own self.
\item No; they support prevention and well-being.
\item Six cleansing actions; trataka cleans the eyes.
\item Kapalabhati, to clear water left in the nose.
\item Bhalabhati.
\item From the top down: neck, shoulders and arms, trunk, hips, knees, ankles and feet.
\item Sukshma vyayama is subtle; sthula vyayama is gross.
\item Steady and comfortable (\emph{sthira-sukham āsanam}).
\item Flexion, extension, lateral flexion, rotation and axial extension.
\item Health (arogya), steadiness (sthairya) and lightness of the body (angalaghava).
\item Puraka, kumbhaka and rechaka.
\item Dharana holds the mind on one object with effort; dhyana is unbroken attention on it.
\end{enumerate}

\paragraph{Sources.} The official guidebook for Level I (2016, older scheme); the Common Yoga Protocol, Ministry of AYUSH; the WHO Constitution (1948); the Yoga Sutra; the Hatha Pradipika (English translation on Wikisource); the Gheranda Samhita (translations found by web search); the Board syllabus for Level 1.
